\section{PLANTEAMIENTO DEL PROBLEMA}
\subsection{Problema Central}
¿Cómo preservar la integridad y mantener la confidencialidad en sistemas de votación electrónica frente a amenazas de seguridad actuales y futuras?




\subsection{Problemas secundarios}
Para el desarrollo del trabajo de investigación se identificaron los siguientes problemas:

\begin{itemize}
	
	
	\item La mayoría de los sistemas de votación electrónica actuales fueron diseñados pensando en amenazas de hoy, sin considerar cómo podrían evolucionar los ataques en los próximos años.

	
	\item No se ha determinado cuánto tiempo adicional tomaría emitir y contar un voto si se usa un cifrado más seguro, lo cual genera dudas sobre su viabilidad práctica.
	
		\item Los sistemas de votación electrónica actuales solo se prueban con herramientas externas básicas, lo cual hace que no se sepa con certeza si resistirían un ataque de una computadora cuántica.
	
	
		\item Los estudios que revisan la seguridad de los sistemas de votación suelen enfocarse en que el voto se cuente correctamente, sin revisar si el cifrado resistiría una computadora cuántica, dejando ambos aspectos sin evaluarse juntos.
	
	
		\item Los desarrolladores de sistemas de votación priorizan la facilidad de uso y la rapidez del sistema, lo cual hace que la seguridad a largo plazo del cifrado quede en segundo plano.
	

\end{itemize}